	\documentclass[12pt]{exam}
	\usepackage[utf8]{inputenc}
	\usepackage[portuguese,brazil]{babel}
	\usepackage[margin=.65in]{geometry}
	\usepackage{amsmath,amssymb}
	\usepackage{multicol}
	\usepackage{multirow}
	\usepackage{float}
	\usepackage{enumerate}
	\usepackage{enumitem}
	\usepackage{array}
	\usepackage{graphicx}
	\usepackage{listings}
	\usepackage{circuitikz}    
	\usepackage{xcolor}
	\usepackage{tikz}
	\usetikzlibrary{shapes.multipart, calc}
	\usepackage{rotating}
	\usepackage{caption}
	%\usepackage[version=3]{mhchem}
	
	\pointpoints{ponto}{pontos}
	\hpword{Pontos:}
	\vpword{Pontos}
	\htword{Total}
	\vtword{Total:}
	\vsword{Resultado}

	\definecolor{azul}{RGB}{0,32,96}
	\definecolor{verde}{RGB}{0,132,124}
	\definecolor{cinza}{RGB}{0,84,132}
	%\definecolor{cinza2}{RGB}{218,220,216}
	\definecolor{cinza2}{RGB}{132,134,136}
	\definecolor{mygreen}{rgb}{0,0.6,0}
	\definecolor{mygray}{rgb}{0.5,0.5,0.5}
	\definecolor{mymauve}{rgb}{0.58,0,0.82}
	\newcommand*{\vttfamily}{%
	\fontencoding{T1}\fontfamily{lmvtt}\selectfont
	}
	\lstdefinestyle{Cstyle}{%
	language=C,numbers=left, keywordstyle=\color{mymauve}\vttfamily,commentstyle=\color{mygreen}\vttfamily,stringstyle=\color{blue},basicstyle=\scriptsize\vttfamily,numberstyle=\tiny\vttfamily,breaklines=true,stepnumber=1,tabsize=2,xleftmargin=\leftmarginii,escapechar=!,captionpos=b,morekeywords={include, printf,NULL},columns=flexible,escapeinside={(^}{^)},
	frame=l,
	framesep=2.5mm,
	framexleftmargin=2.5mm,
	fillcolor=\color{cinza2!50},
	rulecolor=\color{ballblue},
	rulecolor=\color{blue}
	}

	\lstnewenvironment{C}[1][]
	{
	\lstset{style=Cstyle,#1}
	}
	{}


	\newcommand{\class}{Configuração de redes e servidores}
	\newcommand{\profname}{Gilmar Gomes do Nascimento}
	\newcommand{\examnum}{Servidor de arquivos} \newcommand{\examdate}{\the\day/\the\month/\the\year}
	
	\newcommand*{\oldneg}{\mathord{\nao}}
	\pagestyle{headandfoot}
	\runningheadrule
	\firstpageheader{}{}{}
	\runningheader{\class}{\examnum}{\examdate}
	\runningfootrule
	\firstpagefooter{}{}{}
	\runningfooter{}{Página \thepage\ de \numpages}{}
	%\pagestyle{foot}
	%\lfoot{}
	%\cfoot[]{Página \thepage\ de \numpages}
	%\rfoot{}
	%
	%\runningfootrule
	%\boxedpoints %\pointsinmargin
	%\printanswers
	%\SolutionEmphasis{\color{red}}
	%\definecolor{SolutionBoxColor}{gray}{0.8}
	%\renewcommand{\solutiontitle}{\noindent\textbf{Solução:}\enspace}
	%\renewcommand{\solutiontitle}{\noindent\textbf{Solution:}\par\noindent}

	\begin{document}

	\noindent
	\begin{tabular*}{\textwidth}{l @{\extracolsep{\fill}} r @{\extracolsep{6pt}} l}
	\textbf{Disciplina:} \class &  \multicolumn{2}{c}{\multirow{4}{*}{
	\includegraphics[scale=0.35]{logo.png}}}\\
	\textbf{Prof.:} ~\profname \vspace{0.2in} && \\
	\vspace{0.2in} \textbf{Nome:} \makebox[3in]{\hrulefill} \vspace{0.2in} &&\\
	%\textbf{Matrícula:} \makebox[1.5in]{\hrulefill} \textbf{Nota:} \makebox[1in]{\hrulefill} &&\\ 
	%\multicolumn{3}{c}{\textbf{{\examnum \hspace{3px} \examdate}}}\\
	\multicolumn{3}{c}{}\\
	\end{tabular*}\\
	\rule[2ex]{\textwidth}{1.6pt}
	\parskip=0.1in

	%Este exame contêm \numpages\ páginas (incluindo esta capa) e \numquestions\ questões.\\
	%Total de \numpoints~ pontos.

	%\begin{center}
	%Grade Table (for teacher use only)\\
	%\addpoints
	%\gradetable[v][questions]
	%\end{center}

	\noindent
\begin{questions}

\question O que é um SGBD (Sistema Gerenciador de Banco de Dados)? Qual é a sua principal diferença em relação a um servidor de arquivos que vimos na aula anterior?

\begin{solutionbox}{2in} \end{solutionbox}

\question Quais são os SGBDs mais comuns no mercado? Cite pelo menos dois de código aberto e dois proprietários.

\begin{solutionbox}{1.8in} \end{solutionbox}

\question Quais são as portas de rede mais utilizadas pelos SGBDs? Associe cada porta ao SGBD correspondente (ex: MySQL, PostgreSQL, SQL Server).

\begin{solutionbox}{1.5in} \end{solutionbox}

\question O que é o conceito ACID? Explique com suas palavras o que significam as propriedades de Atomicidade, Consistência, Isolamento e Durabilidade.

\begin{solutionbox}{2.5in} \end{solutionbox}

\question Qual é a diferença entre um banco de dados relacional (SQL) e um banco de dados não relacional (NoSQL)? Dê exemplos de cada um.

\begin{solutionbox}{2.5in} \end{solutionbox}

\question O que é DDL (Data Definition Language)? Dê dois exemplos de comandos e explique o que cada um faz.

\begin{solutionbox}{2.2in} \end{solutionbox}

\question O que é DML (Data Manipulation Language)? Dê dois exemplos de comandos e explique o que cada um faz.

\begin{solutionbox}{2.5in} \end{solutionbox}

\question O que é DCL (Data Control Language)? Dê um exemplo de comando e explique sua função.

\begin{solutionbox}{2.5in} \end{solutionbox}

\question Quais são as palavras-chave (comandos) SQL mais utilizadas no dia a dia de quem trabalha com banco de dados? Cite pelo menos cinco.

\begin{solutionbox}{2.5in} \end{solutionbox}

\question O que é um DBA (Administrador de Banco de Dados)? Qual é a sua importância para uma organização?

\begin{solutionbox}{2in} \end{solutionbox}

\question A LGPD (Lei Geral de Proteção de Dados) tem relação com o trabalho de um DBA ou com o armazenamento em bancos de dados? Explique.

\begin{solutionbox}{2.2in} \end{solutionbox}

\question Em quais SGBDs é possível armazenar imagens? Descreva como isso é feito (dica: pense em tipos de dados específicos).

\begin{solutionbox}{2in} \end{solutionbox}

\end{questions}
	\end{document}
